\documentclass[11pt,a4paper]{article}

\usepackage[margin=2.6cm]{geometry}
\usepackage{amsmath,amssymb,amsthm}
\usepackage{graphicx}
\usepackage{booktabs}
\usepackage{microtype}
\usepackage{xcolor}
\usepackage[colorlinks=true,linkcolor=blue!45!black,citecolor=blue!45!black,
            urlcolor=blue!45!black]{hyperref}

\newcommand{\Ns}{N_{s}}
\newcommand{\lI}{\ell_I}
\newcommand{\lR}{\ell_R}
\newcommand{\xj}{x_j}
\newcommand{\muout}{\mu_{\mathrm{out}}}
\newcommand{\dd}{\,\mathrm{d}}

\theoremstyle{plain}
\newtheorem{prop}{Proposition}
\theoremstyle{remark}
\newtheorem*{rem}{Remark}

\title{Spike-interaction asymptotics for the minimal cross-feeding model:\\
how to build the reduced dynamics, and what it says about arrest}
\author{}
\date{\today}

\begin{document}
\maketitle

\begin{abstract}
\noindent
A programme for reducing the three-component cross-feeding PDE to an ODE for
the positions of the population spikes, in the regime $D_N \ll D_I, D_R$. The
reduction has four steps: an outer problem in which each spike is a point sink
of $I$ and a point source of $R$; an inner problem which --- because growth is
linear in $N$ --- is a \emph{linear Schr\"odinger} eigenvalue problem in a well
the spike digs itself, rather than the algebraic balance familiar from
Gray--Scott; a self-consistency condition fixing the spike masses; and a
solvability condition giving the drift law $\dot \xj \simeq x_0^2\,
\partial_x\muout(\xj)$. Every step that can be checked against the arrested
$29$-spike state at $p=1$ is checked here, and the outer problem and the
inner virial identity both hold. The frozen-strength pair force does not: it
predicts an attractive interaction with an $e$-folding time of $7.6\times10^3$
for a run that sat still to $t=10^8$. That discrepancy localises the physics of
arrest in the configuration-dependence of the spike masses, and suggests a
cheap decisive numerical experiment.
\end{abstract}

\section{The model and the state we are trying to explain}

With $k=0$ and unit energy weights the minimal model in one spatial dimension is
\begin{align}
\partial_t N &= D_N N_{xx} + N\,\mu(I,R), &
\mu(I,R) &\equiv (1-l)cI + dR - m, \label{eq:N}\\
\partial_t I &= D_I I_{xx} + K - rI - cNI, \label{eq:I}\\
\partial_t R &= D_R R_{xx} + lcNI - rR - dNR, \label{eq:R}
\end{align}
on a ring of circumference $L$, with $D_R = p$ the scan parameter. Throughout,
numbers come from \texttt{run4\_arrest\_check/state1.jld2}: the arrested
$29$-spike state at $K=15$, $l=0.999$, $m=c=d=r=1$, $D_N=10^{-6}$, $D_I=1$,
$p=1$. Table~\ref{tab:state} collects what that state actually looks like; all
of it is reproduced by \texttt{verify\_spikes.jl} in this directory.

\begin{table}[h]
\centering
\begin{tabular}{lll}
\toprule
quantity & value & \\
\midrule
domain, spikes            & $L=35.4144$, $k=29$          & spacing $d = 1.2212$ \\
spike amplitude, mass     & $A = 255.3$, $S=14.2598$     & identical across spikes to $6$ d.p. \\
spike width               & $x_0 = 0.02204$              & 2nd moment; $= 0.02228$ from $S/A$ \\
compare                   & $\sqrt{D_N/m} = 10^{-3}$     & $x_0/\sqrt{D_N/m} = 22$ \\
$\mu$ at spike core       & $+1.856\times10^{-3}$        & $-0.1443$ at mid-gap \\
$I$                       & $1.0270$ at core             & $3.1630$ at mid-gap \\
$R$                       & $1.00083$ at core            & $0.85254$ at mid-gap \\
resource ranges           & $\lI = \sqrt{D_I/r} = 1$     & $\lR = \sqrt{p/r} = 1$ \\
\bottomrule
\end{tabular}
\caption{The arrested state. Note $R$ is \emph{elevated} at the spike and $I$
\emph{depleted}: $R$ is the activator (self-produced food), $I$ the inhibitor
(the thing competed for).}
\label{tab:state}
\end{table}

Three lengths matter, and they order as
\begin{equation}
  x_0 = 0.022 \;\ll\; d = 1.22 \;\sim\; \lI = \lR = 1 .
\end{equation}
This is the \emph{semi-strong} interaction regime: spikes are narrow, but the
fields that couple them vary on the scale of the spacing itself. In particular
$d/\ell \approx 1.2$, so neighbours are \emph{not} exponentially far apart, and
$e^{-d/\ell}\approx 0.3$ is not small. Whatever makes the coarsening slow here,
it is not the smallness of $e^{-d/\ell}$.

\begin{rem}
$x_0 = 22\sqrt{D_N/m}$. The naive inner scale is wrong by a factor of $22$,
and the reason is structural: \eqref{eq:N} is \emph{linear} in $N$, so there is
no algebraic balance $D_N N'' \sim N^2$ to set a $\operatorname{sech}^2$ width.
The localisation is a linear bound state, and its width is set by the curvature
of the well, which lives on the outer scale. This is the one place where the
calculation genuinely differs from Gray--Scott, and it makes the inner problem
\emph{easier}, not harder.
\end{rem}

\section{Step 1: the outer problem --- spikes as point sources}

Away from the spikes $N$ is negligible ($\sim 10^{-63}$ in the gaps of the
measured state), so \eqref{eq:I}--\eqref{eq:R} linearise:
\begin{equation}
  D_I I_{xx} - rI = -K, \qquad D_R R_{xx} - rR = 0 ,
\label{eq:outer}
\end{equation}
punctured at each $\xj$ by jump conditions obtained by integrating across the
spike. Writing $\sigma^{I}_j, \sigma^{R}_j$ for the signed point-source
strengths,
\begin{equation}
  D_I\big[I_x\big]_{\xj} = -\sigma^I_j = \int_{\text{spike}} \!\!(cNI - K + rI)\dd x,
  \qquad
  D_R\big[R_x\big]_{\xj} = -\sigma^R_j = -\!\!\int_{\text{spike}}\!\! (lcNI - dNR - rR)\dd x .
\label{eq:jumps}
\end{equation}
Measured on the state, over a window of $\pm 0.085$ around a core, both hold to
four figures:
\[
  \sigma^I = -12.926, \qquad \sigma^R = +0.93045, \qquad
  \text{ratio (jump)/(integral)} = 1.0000 \text{ for both.}
\]
The linear terms in \eqref{eq:jumps} are \emph{not} optional: over that window
$\int(K-rI) = 2.44$ against an uptake of $15.36$, so dropping them costs $16\%$.
This is a symptom of $x_0/\ell$ not being very small; see \S\ref{sec:caveats}.

On the ring the solution of \eqref{eq:outer} is a superposition of Green's
functions,
\begin{equation}
  g_\ell(x) = \frac{\ell}{2D}\,
  \frac{\cosh\!\big((L/2 - |x|)/\ell\big)}{\sinh\!\big(L/2\ell\big)},
  \qquad D g_\ell'' - r g_\ell = -\delta(x),
\end{equation}
so that
\begin{equation}
  I(x) = \frac{K}{r} + \sum_k \sigma^I_k\, g_{\lI}(x - x_k),
  \qquad
  R(x) = \sum_k \sigma^R_k\, g_{\lR}(x - x_k).
\label{eq:outerfields}
\end{equation}
This is exact up to the width of the spikes, and it is very accurate: fitting
$I - K/r$ and $R$ over one gap to a single $\cosh$ with the \emph{predicted}
$\lI = \lR = 1$, with only the amplitude free, gives relative residuals of
$4.4\times10^{-4}$ for both fields. The outer half of the reduction is
essentially exact for this state.

\section{Step 2: the inner problem --- a self-dug potential well}

Zoom in on one spike. On the inner scale $I$ and $R$ are quasi-static and
\eqref{eq:N} reads
\begin{equation}
  D_N \Ns'' + \mu(x)\,\Ns = 0 ,
\label{eq:inner}
\end{equation}
a zero-energy Schr\"odinger bound state in the potential $\mu$. Equivalently,
$-D_N\psi'' + V\psi = E\psi$ with $V = -\mu$ and $E = 0$. Two consequences.

\paragraph{The virial identity.} Multiplying \eqref{eq:inner} by $\Ns$ and
integrating over a period,
\begin{equation}
  \boxed{\;D_N \int (\partial_x \Ns)^2 \dd x = \int \mu\,\Ns^2 \dd x\;}
\label{eq:virial}
\end{equation}
exactly, for any spike steady state, whatever its shape. On the measured state
the two sides are $2.5287$ and $2.6253$, a ratio of $0.963$ --- which for a
quantity assembled from a numerical second derivative of a field spanning $65$
orders of magnitude is a clean confirmation. This is the cheapest possible test
that a candidate state really is a steady state of \eqref{eq:N}, and it is
worth adding to the arrest-check battery: it is shape-free, so it cannot be
satisfied by a grid-pinned artefact.

\paragraph{Self-consistency fixes the mass.} Approximating the well
harmonically, $\mu \approx \mu_0 - \tfrac12\kappa\,\xi^2$ with
$\xi = x - \xj$, the ground state is a Gaussian $\Ns \propto
e^{-\xi^2/2x_0^2}$ with
\begin{equation}
  x_0^2 = \sqrt{2D_N/\kappa}, \qquad \mu_0 = \sqrt{D_N\kappa/2}
  \qquad\Longrightarrow\qquad \mu_0\,x_0^2 = D_N .
\label{eq:selfcons}
\end{equation}
The second relation is the self-consistency condition: \emph{the growth rate at
the spike core must sit exactly at the ground-state energy of the spike's own
well}. Since the well is dug by the spike, this is what determines the mass
$S_j$. The corollary $\mu_0 x_0^2 = D_N$ is parameter-free and testable:
measured, $\mu_0 x_0^2 = 9.017\times10^{-7}$ against $D_N = 10^{-6}$, a ratio of
$0.90$. The residual $10\%$ is the non-harmonicity of the well, discussed next.

\paragraph{What the well actually looks like.} The point-source structure of
\eqref{eq:outerfields} means both $I$ and $R$ have a \emph{cusp} at the spike,
so on the inner scale $\mu \approx \mu_0 - \alpha|\xi| - \tfrac12\kappa\xi^2$
with $\alpha = -\big[(1-l)c\,I_x + d\,R_x\big]$ evaluated just outside the
spike. Measured: $\alpha = 0.463$, of which the $R$ part is $-0.4696$ and the
$I$ part only $+0.0065$ --- the tent is dug by the byproduct. If one takes that
tent as sharp, the ground state is an Airy function with
\begin{equation}
  x_0 \sim (D_N/\alpha)^{1/3} = 0.0129, \qquad
  \mu_0 = |a_1'|\,(D_N\alpha^2)^{1/3} = 6.10\times10^{-3},
\end{equation}
against measured $0.0220$ and $1.86\times10^{-3}$: too narrow by $1.7\times$ and
too deep by $3.3\times$. The reason is visible in the numbers --- the source is
not narrow compared with the well it digs, $x_0/(D_N/\alpha)^{1/3} = 1.7$, so
the tent is substantially rounded by the spike's own finite width. The honest
inner problem is therefore the coupled system
\begin{equation}
  D_N \Ns'' + \big[(1-l)cI + dR - m\big]\Ns = 0, \quad
  D_I I'' = cN I, \quad
  D_R R'' = (dR - lcI)N,
\label{eq:innerBVP}
\end{equation}
on $\xi \in (-\infty,\infty)$, with $\Ns \to 0$ and $I,R$ matching the outer
slopes $\sigma^{I,R}/2D$ as $|\xi|\to\infty$. That is a small two-point BVP with
one eigenvalue-like unknown (the mass), solvable by Newton in seconds. It is
worth doing properly rather than harmonically, because it is what supplies the
$O(1)$ constants everywhere downstream.

\section{Step 3: closing the loop --- the $k$-spike quasi-equilibria}

Steps 1 and 2 close on each other. Given positions $\{\xj\}$:
\eqref{eq:outerfields} gives the local field values $(I_j, R_j)$ in terms of the
strengths, \eqref{eq:innerBVP} gives the strengths in terms of $(I_j, R_j)$, and
the fixed point is $k$ algebraic equations for $k$ masses. For equally spaced
spikes symmetry collapses this to a single scalar equation, so the exact
$k$-spike state is available in closed form for every $k$, at the cost of one
Newton solve on one period --- no PDE integration, and no restriction to $k$
that the dynamics happens to visit. This is the object whose Floquet spectrum
gives the stable wavenumber band.

\section{Step 4: motion --- the solvability condition}

Now let the environment be slightly asymmetric. Write $\mu = \mu_s(\xi) +
G_j\,\xi + O(\xi^2)$, where $\mu_s$ is the even part of the well and
\begin{equation}
  G_j \equiv \partial_x \muout(\xj)
   = \sum_{k\neq j}\Big[(1-l)c\,\sigma^I_k\,g_{\lI}'(\xj - x_k)
     + d\,\sigma^R_k\,g_{\lR}'(\xj - x_k)\Big]
\label{eq:G}
\end{equation}
is the gradient of the field due to all \emph{other} spikes. Put
$N = \Ns(x - X(t)) + \phi$ and linearise:
\begin{equation}
  -\dot X\, \Ns' = \mathcal{L}\phi + G\,\xi\,\Ns ,
  \qquad \mathcal{L} \equiv D_N\partial_\xi^2 + \mu_s .
\label{eq:solv}
\end{equation}
Here is the structural difference from Gray--Scott. The kernel of $\mathcal{L}$
is $\Ns$ itself --- \eqref{eq:inner} says so --- and $\Ns$ is \emph{even}.
Projecting \eqref{eq:solv} onto it gives $0 = 0$, because
$\langle \Ns,\Ns'\rangle = 0$ and $\langle \Ns, \xi \Ns\rangle = 0$ by parity.
The usual translation-mode solvability is degenerate, and the drift has to be
extracted one level down, from the odd mode.

Project instead onto the first excited state $\psi_1$, with
$\mathcal{L}\psi_1 = -\omega\psi_1$ and $\omega = \sqrt{2D_N\kappa}$ the
level spacing:
\begin{equation}
  -\dot X\,\langle \psi_1, \Ns'\rangle
  = -\omega\,\langle\psi_1,\phi\rangle + G\,\langle \psi_1, \xi \Ns\rangle .
\end{equation}
Defining $X$ as the centroid removes the dipole content of $\phi$, so
$\langle\psi_1,\phi\rangle = 0$, and for the Gaussian ground state
$\psi_1 \propto \xi\Ns$ the two remaining brackets are $-\sqrt{\pi}/2$ and
$x_0^2\sqrt{\pi}/2$, giving
\begin{equation}
  \boxed{\;\dot \xj \;=\; C\,x_0^2\, G_j, \qquad C = O(1),
  \qquad x_0^2 = \sqrt{2D_N/\kappa} \;\propto\; \sqrt{D_N}\;}
\label{eq:drift}
\end{equation}
with $C = 1$ in the harmonic/Gaussian reduction. The measured mobility is
$x_0^2 = 4.86\times10^{-4}$.

Two things to take from \eqref{eq:drift}. First, the spike climbs the gradient
of its environment: it grows faster on the side where $\mu$ is larger, and the
centroid follows. Second --- and this is the key point for the arrest
question --- the mobility is $O(\sqrt{D_N})$, not $O(D_N)$ and \emph{not}
exponentially small in anything. Combined with $G$ being an $O(1)$ function of
$d/\ell$ rather than an exponentially small one, the reduced dynamics has no
small parameter left with which to make coarsening take longer than $10^8$.

\begin{rem}
The step from \eqref{eq:solv} to \eqref{eq:drift} is a two-mode reduction, not a
systematic expansion. A systematic version has to work with the full linearised
operator including the feedback of $\phi$ on $\mu$ through
\eqref{eq:I}--\eqref{eq:R}, whose kernel \emph{is} $\Ns'$ by translation
invariance and which is not self-adjoint, so the projection is onto its adjoint
null mode. That calculation could change $C$, and in principle could reveal a
mobility much smaller than $x_0^2$; it is the main open point in this document.
\end{rem}

\section{What this predicts, and where it collides with the data}
\label{sec:collide}

\subsection{Sign and range of the interaction}

From \eqref{eq:outerfields} the field of one neighbour contributes to $\mu$ at
weight $(1-l)c\,\sigma^I$ through the $I$ channel and $d\,\sigma^R$ through the
$R$ channel. With the measured strengths,
\[
  (1-l)c\,\sigma^I = -0.01293 \quad(\text{range } \lI = 1), \qquad
  d\,\sigma^R = +0.9304 \quad(\text{range } \lR = \sqrt{p}).
\]
So $I$ mediates \textbf{repulsion} at range $\lI$ (moving towards a neighbour
means moving into its depletion hole) and $R$ mediates \textbf{attraction} at
range $\lR$ (moving towards a neighbour means moving into its byproduct plume).
This is the physical content of the $p$ scan: $p$ sets $\lR/\lI = \sqrt p$, i.e.
whether the attraction is shorter-ranged than the repulsion.

Balancing the two channels, $\;(1-l)c|\sigma^I|e^{-s}/2 = \sigma^R
e^{-s/\sqrt p}/2\sqrt p$, gives a crossover distance
\begin{equation}
  s^{*} = \frac{\sqrt p}{1-\sqrt p}\,
          \log\!\frac{\sigma^R}{(1-l)c|\sigma^I|\sqrt p},
\label{eq:crossover}
\end{equation}
inside which the net force is attractive and outside which it is repulsive ---
so a lattice with $d > s^*$ is stable and one with $d < s^*$ coarsens, and
$s^*(p)$ \emph{is} the predicted arrest spacing. Holding the $p=1$ strengths
fixed (so this is the shape of the $p$ dependence, not a quantitative
prediction away from $p=1$), $\sigma^R/((1-l)c|\sigma^I|) = 72$ and

\begin{center}
\begin{tabular}{lccccc}
\toprule
$p$      & $0.01$ & $0.03$ & $0.1$ & $0.3$ & $1$ \\
$s^*$    & $0.731$ & $1.263$ & $2.510$ & $5.908$ & --- (no crossover) \\
\bottomrule
\end{tabular}
\end{center}

\noindent
At $p=1$ the ranges coincide, the $R$ channel wins by a factor of $72$ at every
distance, and there is no stable spacing at all.

\subsection{The collision}

That last statement is contradicted by the data, and not marginally. Take the
arrested state, displace one spike by $\delta$ towards its right neighbour, and
evaluate \eqref{eq:G} and \eqref{eq:drift} with the strengths held fixed:
\begin{center}
\begin{tabular}{lccc}
\toprule
$\delta$ & $G$ & $v = x_0^2 G$ & $e$-folding time $\delta/v$ \\
\midrule
$0.0122$ ($1\%$ of $d$)  & $+3.30\times10^{-3}$ & $+1.61\times10^{-6}$ & $7.6\times10^{3}$ \\
$0.1221$ ($10\%$ of $d$) & $+3.31\times10^{-2}$ & $+1.61\times10^{-5}$ & $7.6\times10^{3}$ \\
\bottomrule
\end{tabular}
\end{center}
The lattice is predicted to be linearly unstable with growth rate
$1.3\times10^{-4}$. The run reached $t = 10^8$ without moving: that is
$1.3\times10^{4}$ $e$-folding times. The frozen-strength reduction is wrong by
an enormous margin, and the interesting question is which of its ingredients
fails. The candidates, in the order I would test them:

\begin{enumerate}
\item \textbf{The strengths are not frozen} (most likely). $\sigma^I_k$ and
      $\sigma^R_k$ are outputs of the self-consistency in \S3, and they depend on
      the whole configuration. When spike $j$ moves right, its two neighbours
      acquire \emph{different} masses, and the resulting asymmetry in
      \eqref{eq:G} is a force channel that the frozen calculation discards
      entirely. In Gierer--Meinhardt and Gray--Scott this is exactly the channel
      that supplies the stabilising repulsion. Concretely: the positions and the
      masses must be evolved as one coupled system, not positions in a fixed
      landscape.
\item \textbf{The mobility is much smaller than $x_0^2$}, because the systematic
      version of \S4 (see the remark) produces a different projection. This
      would need roughly four orders of magnitude, which seems a lot, but the
      degeneracy of the leading solvability condition is exactly the kind of
      structure that produces anomalously small mobilities.
\item \textbf{The drift is real but sub-threshold in the diagnostics.} This one
      is already close to excluded: $v \sim 10^{-5}$ is nine orders above the
      $\sim 10^{-14}$ noise floor of \texttt{peak\_drift}, and the run would have
      merged thousands of times over.
\end{enumerate}

\section{The other eigenvalue family}

Position modes are only half the spectrum. Perturbing the \emph{masses},
$S_j \to S_j + \epsilon\, e^{\lambda t + i q j}$, and solving the outer response
produces a nonlocal term $\propto \int \phi$ in the inner eigenvalue problem ---
a nonlocal eigenvalue problem (NLEP) of the standard type. Its $q \neq 0$ branch
is the competition instability: one spike grows at its neighbour's expense and
the neighbour dies. This family is
\begin{itemize}
\item $O(1)$ in $D_N$ rather than $O(\sqrt{D_N})$, so it acts first and fast;
\item the natural explanation for a \emph{box-independent} arrest spacing, which
      is what run1 ($10$ peaks) and run2 ($100$ peaks) show across a $10\times$
      change of $L$;
\item not obtainable from the position asymptotics at all.
\end{itemize}
If the arrest is really set by the competition threshold $k_c$, then the
positions can be marginal or even weakly unstable without contradicting the
observations --- provided the position instability is slow enough, which
\S\ref{sec:collide} says it is not, on the frozen-strength estimate. Resolving
that tension is the crux.

\section{A numerical programme, cheapest first}

\begin{enumerate}
\item \textbf{Displace one spike and measure the drift.} Shift a single spike in
      the arrested state by $10\%$ of the spacing, re-relax $I$ and $R$ (or just
      evaluate the residual), and run \texttt{peak\_drift}. The prediction is
      $v = +1.6\times10^{-5}$ towards the near neighbour. The \emph{sign} says
      attraction or repulsion, i.e.\ unstable or stable lattice; the
      \emph{magnitude} tests \eqref{eq:drift}. Both $\mu_x$ and $x_0$ are
      measurable directly off the perturbed state, so this tests the drift law
      with no asymptotics required to evaluate it. This is an afternoon's work
      and it discriminates between the three candidate failures above.
\item \textbf{Add the virial check \eqref{eq:virial} to the arrest battery.}
      Shape-free, cheap, and it cannot be faked by grid pinning.
\item \textbf{Solve the inner BVP \eqref{eq:innerBVP} by Newton.} Gives the
      $O(1)$ constants and a spike profile to compare with the PDE state.
\item \textbf{Build the $k$-spike quasi-equilibria (\S3) and their Floquet
      spectrum.} Both eigenvalue families at once, for every $k$, without
      integrating the PDE. This is the object that answers ``is the arrest
      spacing at the edge of the stable band''.
\item \textbf{Only then} the analytic NLEP, if the numerics says the competition
      threshold is what matters.
\end{enumerate}

\section{Caveats}
\label{sec:caveats}

\begin{itemize}
\item \textbf{$x_0/\ell$ is not very small.} $x_0/\ell = 0.022$ sounds
      comfortable, but the effective expansion parameter keeps showing up as
      something more like $x_0/(D_N/\alpha)^{1/3} = 1.7$ or the $16\%$ correction
      in \eqref{eq:jumps}. Expect $O(1)$ constants to be genuinely $O(1)$, and
      do not trust a factor of two from this calculation.
\item \textbf{$l = 0.999$ is a near-singular limit.} $(1-l)cI \approx 10^{-3}$ is
      negligible in the \emph{value} of $\mu$ but its gradient,
      $(1-l)c\,\Delta I \approx 2\times10^{-3}$, is the same size as $\mu_0$
      itself. The $I$ channel cannot be dropped as small; it is small in value
      and $O(1)$ in effect.
\item \textbf{No variational structure.} There is no free energy running
      downhill here, so nothing forces coarsening to proceed, and a stable
      multi-spike state is not anomalous. The intuition that it must coarsen is
      imported from Cahn--Hilliard and does not apply.
\item \textbf{Quasi-staticity of $I,R$ is safe.} The resources relax at $O(1)$
      rates against a drift of $O(10^{-5})$, so the adiabatic elimination in \S1
      is very well satisfied --- five orders of margin.
\end{itemize}

\section*{Literature}

The closest templates: Doelman, Gardner \& Kaper and Kolokolnikov, Ward \& Wei
for semi-strong interaction and competition/oscillation thresholds in
Gray--Scott; Iron, Ward \& Wei for the Gierer--Meinhardt spike dynamics and the
NLEP machinery; Ei, \emph{The motion of weakly interacting pulses}, for the
projection method behind \S4; and Sherratt's Busse-balloon computations for the
Klausmeier model, which is the closest structural analogue --- slowly diffusing
organism plus fast resource, banded patterns, arrested coarsening --- and where
the stable-wavenumber-band programme has been carried out end to end.

\bigskip
\noindent\emph{All numbers quoted here are produced by} \texttt{verify\_spikes.jl}
\emph{in this directory; run it from the repository root with}
\texttt{julia --project=. rest/more\_ai\_derivations/verify\_spikes.jl}.

\end{document}
